\chapter{利息：时间是有价格的}
\chaptermeta{03}

\begin{ledetext}
银行靠利差活着。可利息本身是什么？为什么钱借出去就该收回更多？答案不在钱上，在时间上。
\end{ledetext}

\section{亚里士多德算错了}

亚里士多德对利息意见很大。他说钱是死的，不会像牛羊那样下崽，靠钱生钱违反自然。

他算错的地方在于，利息本来就不是钱在下崽。

利息是你为"提前拿到"付的钱。拆开看，任何一个利率里都塞着几样东西。

\begin{flowlist}
  \flowitem{01}{时间偏好，也就是人天生不耐}{现在给赵磊 8000 块，和一年后给他 8000 块，他选现在。不是因为他算过通胀，是因为人就这样：现在的确定，胜过将来的可能。要他愿意等，得加钱。这笔加钱是时间本身的价格。}
  \flowitem{02}{通胀预期}{一年后那 8000 块买的东西大概率少一点。这部分补的不是等待，是购买力缩水。没人愿意借出一头牛，一年后收回一只羊。}
  \flowitem{03}{违约风险}{借给财政部，和借给赵磊他姐在南京开的那家卤味店，不可能一个价。店可能明年就关了。多收的那部分，是为"可能收不回来"买单。}
  \flowitem{04}{流动性溢价}{三年定期比活期利率高，多出来的部分补的是这三年你不能动它。急用钱的时候，钱在手上和钱在合同里是两样东西。}
\end{flowlist}

赵磊他姐去年借了 41 万经营贷，年化 6.5\%。这 6.5\% 是一层一层摞上去的。

\begin{tablewrap}
\begin{xltabular}{\linewidth}{>{\raggedright\arraybackslash}p{0.20\linewidth}Xrr}
\toprule
\thead{一层一层往上加} & \thead{这一层在补偿什么} & \thead{加多少} & \thead{累计} \\
\midrule
\textbf{纯粹的不耐} & 人就是不愿意等 & +0.5\% & 0.5\% \\
\textbf{通胀补偿} & 一年后同样的钱买得更少 & +1.0\% & 1.5\% \\
\textbf{\textbf{= 名义无风险利率}} & 对照：7 天逆回购 1.4\% & — & \textbf{约 1.5\%} \\
\textbf{流动性溢价} & 这笔钱三年抽不回来 & +1.0\% & 2.5\% \\
\textbf{违约风险溢价} & 一百家小店里总有几家关门 & +4.0\% & \textbf{约 6.5\%} \\
\bottomrule
\end{xltabular}
\end{tablewrap}

这是示意性拆解，不是任何机构的真实报价。风险怎么定价，第 13 章专门讲。

利率从来不是一个数字，是一叠补偿的总和。看到 6.5\% 就骂高利贷，跟看到 1.5\% 就说钱不值钱，犯的是同一个错：只看总和，不看结构。

\section{你妈的定期}

合同上写的那个数叫\term{名义利率}{nominal rate}\index{039@名义利率}。\term{实际利率}{real rate}\index{051@实际利率}约等于名义利率减通胀。这叫费雪关系。

\begin{egbox}{举个栗子}
赵磊他妈在盐城，把 11.6 万存了一年定期，利率 2\%。到期账上躺着 118,320 元。她挺满意，说你看钱变多了。

如果这一年通胀是 3\%，去年花 11.6 万能买到的那一篮子东西，今年要 119,480 元。

钱多了 2320 块，能买的东西少了 1100 出头。实际收益率是 −0.97\%。

银行没骗她。账单上每一分利息都是真的。

偷东西的那位不开对账单。

\end{egbox}

以后看到任何一个收益率，先在心里做一次减法。名义数字是给你看的，实际数字才是你赚的。

\section{这本书最重要的一个除法}

年利率 5\%，今天的 95.24 元一年后变成 100 元。

反过来读：一年后的 100 元，今天只值 95.24 元。

这个折回今天的动作叫\term{贴现}{discounting}\index{059@贴现}，那个 5\% 叫\term{贴现率}{discount rate}\index{060@贴现率}，算出来的 95.24 叫\term{现值}{present value}\index{068@现值}。

把时间拉长，威力就出来了。

\begin{tablewrap}
\begin{xltabular}{\linewidth}{>{\raggedright\arraybackslash}p{0.20\linewidth}Xr}
\toprule
\thead{未来的 100 元} & \thead{怎么折} & \thead{今天值多少} \\
\midrule
\textbf{1 年后} & 100 ÷ 1.05 & 95.24 元 \\
\textbf{10 年后} & 100 ÷ 1.05\textsuperscript{10} & 61.39 元 \\
\textbf{30 年后} & 100 ÷ 1.05\textsuperscript{30} & 23.14 元 \\
\bottomrule
\end{xltabular}
\end{tablewrap}

世界上几乎所有金融资产（股票、债券、房子、一家公司、一条收费公路）本质上都是一串未来的现金流。给它们定价，就是把这串现金流一笔一笔折回今天，再加起来。

分母上坐着的那个贴现率，来自利率。

\begin{pullquote}{}
利率一动，所有资产的价格都要重算。\textbf{没有例外。}
\end{pullquote}

\begin{egbox}{举个栗子：分母变一下，价格砍一半}
假设有项资产每年稳定给你 3.8 万，永远给下去。它值多少钱？

价格等于每年现金流除以贴现率。

贴现率 2\%：3.8 万 ÷ 0.02 = \textbf{190 万}。\\ 贴现率 4\%：3.8 万 ÷ 0.04 = \textbf{95 万}。

现金流一分没少，经营一切照旧，\hlmark{价格砍掉一半}。只因为分母变了。

2022 年以来全球资产价格的剧烈波动，核心就是这个机制。美联储把政策利率从接近零抬到 5\% 以上，全球股债一起跌。跌得最狠的是那些现金流全在很远的将来的资产：不赚钱的成长股、长久期债券。分母抬高对遥远的现金流最致命，30 年后的 100 元本来就只值 23 块，除数一大就不剩什么了。

上一章那道题的答案也在这儿。一张 30 年期、票息 3\% 的债券，市场利率上行 1 个百分点，价格大约跌 20\%。硅谷银行资产端那 150 亿美元浮亏就是这么来的。它不是被人骗了，它是被一个除法干掉的。

\end{egbox}

\section{72 除以 3.5}

贴现是把未来折回今天。反过来把今天推到未来，就是\term{复利}{compound interest}\index{018@复利}：利息也开始生利息。

有个心算工具叫 72 法则。72 除以年化收益率，得到本金翻倍要多少年。

72 ÷ 3.5 ≈ 20.6 年。

精确算是 20.1 年。72 法则在 6\% 到 10\% 之间最准，利率低的时候会略微高估。这点误差不影响它当心算工具，但它提醒了一件事：所有关于复利的漂亮说法，对收益率假设都极其敏感。

赵磊手上有 4.6 万闲钱。按 3.5\% 滚 30 年是 129,122 元，按 2\% 滚 30 年是 83,306 元，差 45,816 元。

而 3.5\% 和 2\% 在直觉上是"差不多的小数字"。

复利的真相和关于复利的谎言，第 22 章掰扯。

\section{两个百分点，128 万}

赵磊今年 32 岁，在南京做产品经理，去年买房贷了 300 万，30 年，等额本息。

\begin{tablewrap}
\begin{xltabular}{\linewidth}{>{\raggedright\arraybackslash}p{0.20\linewidth}rrr}
\toprule
\thead{贷款条件} & \thead{月供} & \thead{30 年总还款} & \thead{其中利息} \\
\midrule
\textbf{300 万 / 30 年 / 年利率 \textbf{3.5\%}} & 约 13,469 元 & 约 484.9 万 & 约 \textbf{184.9 万} \\
\textbf{300 万 / 30 年 / 年利率 \textbf{5.5\%}} & 约 17,036 元 & 约 613.3 万 & 约 \textbf{313.3 万} \\
\textbf{差额} & 每月多 3,567 元 & 多还 128.4 万 & 多付 \textbf{128.4 万} \\
\bottomrule
\end{xltabular}
\end{tablewrap}

盯住最后一行。

纸面上只差 2 个百分点，听起来像四舍五入就没了的小数。换成现实，是赵磊每个月多掏 3567 块，三十年多付出去 128.4 万。在南京这是一辆不错的车加一整套装修，在他妈住的盐城，这是一套小房子。

还有个数更刺眼：3.5\% 那一栏，他借 300 万，最后要还 485 万。5.5\% 那一栏，他还的钱是借的两倍多。

他买的不只是房子。他还买了三十年的等待权，而这个等待权明码标价。

\begin{warnbox}{小心}
看完这张表，很多人第一反应是赶紧提前还贷。这个决定取决于你的贷款利率、你能拿到的其他资产回报、你的现金流稳定性和应急储备，是一道个人化的比较题，没有标准答案。\textbf{本书不提供投资或借贷建议}，只负责让你看懂这张表在说什么。

\end{warnbox}

\section{全世界的宗教都恨过它}

收息现在天经地义。它曾经是重罪。

《旧约》禁止向同胞取利。中世纪教会把放贷取息定为罪，1179 年第三次拉特兰会议规定放高利贷者不得领圣餐，死后不许葬进教会墓地。伊斯兰教法至今禁 riba，所以伊斯兰金融要用买卖、租赁、合伙分成的结构绕开固定利息。

反应这么激烈是有原因的。

古代社会里借钱的人几乎都是走投无路的穷人。家里歉收，孩子快饿死了，才去借粮。这时候有人在你身上加三成，你还不上就拿女儿抵债。这不是金融，是趁火打劫。禁令针对的是这种\term{消费性借贷}{consumption loan}\index{071@消费性借贷}。

转折发生在借钱的用途变了以后。

商人借钱买一船货，工厂主借钱买一台织机，借来的钱产生了新的产出。利息的性质跟着变了：它不再是从穷人碗里抢饭，是出资人分掉自己出资创造出来的那部分增量。

现代社会接受的是这套逻辑：资本有机会成本，风险该被补偿，钱流向回报更高的地方有效率。

那条古老的道德直觉并没有消失。今天各国给消费贷设利率上限、打击非法放贷、强制明示年化利率，护的还是那个走投无路才借钱的人。

\section{从倒贴，到 1.0\%}

利息为正听起来天经地义。人类历史上绝大多数时间确实如此。

2014 年，事情变得很奇怪。

那一年欧洲央行把存款便利利率降到 −0.1\%。2016 年日本银行跟进。丹麦、瑞士、瑞典也进过负值区间。

负利率的字面意思是：你借钱给别人，还要倒贴保管费。把 100 万借给德国政府，十年后拿回 99 万。2020 年，全球负收益率债券的规模一度超过 17 万亿美元。

这在人类金融史上极其罕见。它是通缩、老龄化和长期需求不足逼出来的：央行想刺激经济，发现利率已经到零，只好把它按到水面以下。

这场实验现在结束了。日本这两年走出来的路径，是最值得看的一条线。

\begin{flowlist}
  \flowitem{01}{2024 年 3 月}{结束负利率与收益率曲线控制，政策利率从 −0.1\% 抬到 0\textasciitilde{}0.1\%。十七年来第一次加息。}
  \flowitem{02}{2024 年 7 月}{加到 \textbf{0.25\%}。市场剧烈震荡，套息交易开始被迫平仓。}
  \flowitem{03}{2025 年 1 月}{加到 \textbf{0.5\%}。}
  \flowitem{04}{2025 年 12 月}{加到 \textbf{0.75\%}。}
  \flowitem{05}{2026 年 6 月 16 日}{加到 \textbf{1.0\%}，\hlmark{1995 年 9 月以来的最高水平}。底气来自工资和物价：2026 年"春斗"薪资涨幅预期约 5.02\%，5 月 PPI 同比 6.3\%。}
\end{flowlist}

欧洲央行也在 2026 年 6 月 11 日上调了存款便利利率，三年来第一次加息。

\begin{statrow}{4}
  \statitem{3.50}{–3.75\%}{美国联邦基金利率目标区间}
  \statitem{3.0}{\%}{中国 1 年期 LPR}
  \statitem{3.5}{\%}{中国 5 年期以上 LPR}
  \statitem{1.0}{\%}{日本政策利率}
\end{statrow}

\begin{databox}{2026 数据（截至 2026 年 8 月）}
美联储的联邦基金利率目标区间维持在 \textbf{3.50\%–3.75\%}，年内已连续 5 次按兵不动，理由是通胀仍然高企。美国 10 年期国债收益率在 8 月底升至 2025 年 1 月以来的最高水平。

中国方面，8 月 20 日 LPR 报价 1 年期 \textbf{3.0\%}、5 年期以上 \textbf{3.5\%}，均已连续多月持平；7 天期逆回购政策利率 \textbf{1.4\%}；货币政策基调是"适度宽松"。日本政策利率 \textbf{1.0\%}。

\end{databox}

同一天，同一个地球，钱的价格从 1.0\% 到 3.75\%，差了将近三倍。

这个差价不会安静待着。在日本借 1\% 的钱，换成美元买 3.5\% 的资产，中间的利差是白拿的，只要汇率别动。这叫\term{套息交易}{carry trade}\index{058@套息交易}，是全球市场最大的一条暗流，也是 2024 年 8 月那场全球股市急跌的导火索。它怎么转、为什么会突然反转，第 12 章算这笔账。

\begin{mythbox}{常见误解}
\mvfrow{误解}{"利率当然是越低越好，借钱便宜，大家都开心。"}
\mvfrow{事实}{这个说法是错的。利率是经济的体温，长期贴地通常说明病人虚弱，不是身体好。它意味着需求不足、投资回报低到央行必须一直踩着油门。\\ 代价也是实的：低利率鼓励人借钱去买存量资产（房子、股票）而不是投新产能，把资产价格推上去，对有房有股的人是福音，对想上车的年轻人是灾难；它让本该出清的僵尸企业靠低息续命；它压垮靠利息生活的退休者，以及必须用利息覆盖长期给付的养老金和保险公司。日本和欧洲用了整整十年才从这个坑里爬出来。\\ 至于低利率到底多大程度上推高了资产价格、多大程度上只是反映了增长本身在放缓，学界吵了十几年，目前没有共识。}
\end{mythbox}

\begin{mythbox}{常见误解}
\mvfrow{误解}{"存款利率一降再降，是银行黑心，赚我们的钱。"}
\mvfrow{事实}{因果方向反了。2026 年二季度末，中国商业银行净息差是 \textbf{1.41\%}，历史低位。银行两头被压：资产端的贷款利率跟着 LPR 和 1.4\% 的政策利率往下走，负债端的存款利率不跟着降，利差会直接变负。\\ 存款利率是被贷款利率拖着走的，而贷款利率反映的是整个经济的资金供求和政策取向。这不是说银行是慈善机构，是说这件事的因果链条，和直觉指的方向正好相反。传导的完整机制，第 9 章讲。}
\end{mythbox}

\bookbreak

学会了贴现，后面所有章节的门都开了。债券为什么涨跌、股票为什么被杀估值、房子值不值这个价、一家公司该值多少钱，全是同一个除法。

赵磊的房贷合同上写着一个年利率。

这个数字里有多少是给时间的，多少是给通胀的，多少是给"他可能失业"这件事的，合同上不写。

\begin{takeaway}{本章一句话}
利息不是钱自己生钱，是时间偏好、通胀预期、违约风险和流动性溢价叠出来的价格。它反过来用就是贴现率，而贴现率一动，世界上每一项资产都要重新标价。

\begin{itemize}
  \item 实际利率约等于名义利率减通胀。存款 2\%、通胀 3\%，一年白送出去 1100 块。
  \item 5\% 贴现率下，30 年后的 100 元今天只值 23.14 元。
  \item 300 万房贷 30 年，3.5\% 和 5.5\% 的利息总额差 128.4 万。
  \item 2026 年 8 月：美国 3.50\%–3.75\%、中国 LPR 3.0\%/3.5\%、日本 1.0\%。利差就是资金流动的燃料。
\end{itemize}

\end{takeaway}

\begin{bridgebox}
\textbf{下一章}利率里有一块叫"通胀补偿"。那通胀到底是什么？它凭什么每年从你口袋里拿走一点，还让你以为是自己不会理财？
\end{bridgebox}

